Put $Q=H(\psi)$, $S=\operatorname{tr}Q^{-1}$ and
$Lw=\operatorname{Re}\operatorname{tr}(Q^{-1}H(w))$. The rank-one Hessian rule of Lemma~\ref{ms:lemma-3-1} shows that the real principal symbol is $\widehat\xi^*Q^{-1}\widehat\xi$ for the four-slot encoding $\widehat\xi$ of a real covector. The comparison principle follows by integrating the derivative of $\log\det$ along the straight segment of two positive Hessians. The averaged operator has positive definite real symbol on the compact closure, so the maximum principle applies. The formulas for differentiating $\log\det$ and its concavity are those of Lemma~\ref{ms:lemma-3-2}.

Both endpoint densities are strictly above $\det H(B_0)$, with a uniform positive gap by compactness. Choose $0<b'<1<b$ so that
\[
 \det H(B_0+b'r)<g_t<\det H(B_0+br)
 \quad\text{on }\overline{\widetilde\Omega},\quad 0\le t\le1.
\]
The lower inequality follows from the gap by taking $b'$ small, and the upper one from $H(r)\ge cI$ by taking $b$ large. The comparison principle gives $br\le u\le b'r$. Smooth division in a boundary collar gives $u=\alpha_0r$ with $b'\le\alpha_0\le b$, so $\nabla u=\alpha_0\nabla r$ on the boundary. Also $Lr\ge cS$, and the arithmetic-geometric mean inequality gives $S\ge mg_t^{-1/m}\ge c_0>0$. For any constant real unit direction $e$, differentiating once and using concavity on differentiating twice gives
\[
 L\psi_e=(\log g_t)_e,\qquad L\psi_{ee}\ge(\log g_t)_{ee}.
\]
The maximum principle for $\pm\psi_e+Cr$ bounds all gradients from their boundary values, already bounded by the barriers and the data. Applied to $\psi_{ee}+Cr$, it reduces the upper real Hessian bound to a boundary estimate.

\emph{Boundary second derivatives.} In a flat chart centered at a boundary point let $\nu$ be the outward real unit normal, identified with its quaternionic column. Tangential real directions $v,w$ satisfy $D^2u(v,w)=\alpha_0D^2r(v,w)$ at the center. For a mixed derivative fix a real unit tangent $e$ and construct a local quaternionic rigid-motion field $X=(e+Ax)\cdot\nabla$ by setting
\[
 v=\frac{P_{\rm tan}D^2r\,e}{|\nabla r|},\qquad
 A=v\nu^*-\nu v^*-\nu(\nu^*v)\nu^* .
\]
Here $P_{\rm tan}$ is the real tangent projection. As $\operatorname{Re}(\nu^*v)=0$, one has $A^*=-A$ and $A\nu=v$. At the center $Xr=0$ and for every real tangent $w$,
\[
 \partial_w(Xr)=D^2r(w,e)-|\nabla r|\,w\cdot A\nu=0.
\]
In particular $|Xu|=|\alpha_0Xr|\le C|x|^2$ along the boundary in the chart. The rigid-motion flow acts by quaternionic unitary rotations and translations. Covariance gives $L(X\psi)=X\log g_t$: any commutator from conjugating $Q$ pairs to zero with $Q^{-1}$. Subtract $L(XB_0)$ to get $|L(Xu)|\le CS$, using the smooth data and $S\ge c_0$.

In a small fixed chart ball take $w_0=C_1|x|^2-C_2r$. First choose $C_1$ large so $w_0\ge|Xu|$ on the true boundary by the quadratic bound and on the spherical boundary by the gradient bound. Then take $C_2$ large enough for $L(w_0\pm Xu)\le0$, using $H(|x|^2)=8I$ and $Lr\ge cS$. The minimum principle gives $|Xu|\le w_0$. Both vanish at the center, so the bound controls $|\partial_\nu Xu|$ there. This is $|D^2u(\nu,e)|$, because the extra term $\nabla u\cdot A\nu$ vanishes with $\nabla u$ normal and $A\nu$ tangent. Compactness makes the estimates uniform over the boundary.

By the product rule of Lemma~\ref{ms:lemma-3-1}, at the boundary
\[
 H(u)=\alpha_0H(r)+G(\alpha_0)G(r)^*+G(r)G(\alpha_0)^*,
 \qquad G(r)=|\nabla r|\nu .
\]
The cross terms vanish on the quaternionic orthogonal complement of $\nu$. Thus $Q$ compressed to $\nu^\perp$ is bounded below by $b'H(r)|_{\nu^\perp}\ge cI$. In a quaternionic orthonormal basis beginning with $\nu$, the real rank-one Hessian rule therefore gives
\[
 Q=\begin{pmatrix}\psi_{\nu\nu}+c_1&z^*\\z&D\end{pmatrix},
\]
where $c_1,z,D$ depend only on the already bounded real tangential or mixed second derivatives. For $m>1$, Lemma~\ref{ms:lemma-3-2} gives
\[
 \psi_{\nu\nu}=g_t/\det D+z^*D^{-1}z-c_1,
\]
which is bounded because $D\ge cI$. For $m=1$ use $Q=g_t$. The upper Hessian estimate in the interior now gives $D^2\psi\le CI$. Since $\Delta\psi=\operatorname{tr}Q>0$, every real Hessian eigenvalue is bounded below by $-(4m-1)C$. This proves a uniform full $C^2$ bound.
